\section{Ultraviolet--infrared hierarchy}

\subsection{UVIR-001: the minimal condensate route closes negative}

The tested microscopic candidate is
\begin{equation}
 S_\Phi=\int \dd^4x\sqrt{-g}\left[
 -\nabla_\mu\Phi^*\nabla^\mu\Phi
 -m^2|\Phi|^2
 -\frac{\lambda_4}{2}|\Phi|^4
 -\frac{\lambda_6}{3\Lambda^2}|\Phi|^6\right].
 \label{eq:candidate-uv}
\end{equation}
Writing $\Phi=\rho e^{i\Theta}/\sqrt2$ and
$Z=-\nabla_\mu\Theta\nabla^\mu\Theta$, its stable finite-density branch can
be reduced to a phase action $P(Z)$.  In the pure-sextic limit the explicit
amplitude integration gives
\begin{equation}
 P(Z)=\frac{2\Lambda}{3\sqrt{\lambda_6}}(Z-m^2)^{3/2}.
 \label{eq:sextic-pz}
\end{equation}
This nonanalytic power does not supply the required static spatial operator.
About a timelike background $Z_0=\mu^2$, a static phase gradient
$q=|\nabla\vartheta|$ instead yields
\begin{equation}
 P(\mu^2-q^2)-P(\mu^2)=-P_Z(\mu^2)q^2+\mathcal O(q^4),
 \qquad P_Z=\frac{\rho_0^2}{2}>0.
 \label{eq:uvir001-static}
\end{equation}
The leading response is therefore quadratic for every stable nonzero branch
with the canonical phase kinetic term.  UVIR-001 is closed negative for this
minimal candidate.  The result rejects neither the condensate's density and
topology roles nor a bottom-up spatial $Y^{3/2}$ EFT.

\subsection{UVIR-002: provisional two-sector route}

UVIR-002 also rejects a standalone nonrelativistic gradient-dominated branch:
the branch producing the desired cubic spatial action has a negative quadratic
time kinetic coefficient.  The least-incomplete route tested separates the
jobs of the condensate and force mediator.  The independent invariants are
\begin{equation}
 \mathcal Q=U^\mu\nabla_\mu\psi,
 \qquad
 \mathcal Y=h^{\mu\nu}\nabla_\mu\psi\nabla_\nu\psi.
 \label{eq:preferred-invariants}
\end{equation}
Here $\Phi$ retains finite density, circulation, defects and topology, while
$(\psi,U^\mu)$ carries the conditional low-acceleration response.  This is a
provisional architecture selection, not microscopic closure.

\subsection{UVIR-003 Stage A action and limited pass}

Stage A chooses an independent unit vector and retains the complete
Einstein--aether two-derivative basis,
\begin{align}
 \mathcal L_U=-\frac{M_U^2}{2}\big[&c_1(\nabla_\mu U_\nu)(\nabla^\mu U^\nu)
 +c_2(\nabla_\mu U^\mu)^2 \nonumber\\
 &+c_3(\nabla_\mu U_\nu)(\nabla^\nu U^\mu)-c_4a_\mu a^\mu\big]
 +\frac{\lambda}{2}(U^\mu U_\mu+1),
 \label{eq:aether-action}
\end{align}
where $a_\mu=U^\nu\nabla_\nu U_\mu$.  A positive alignment energy couples
the frame to the condensate current without imposing equality.  The regulated
force truncation is
\begin{equation}
 \mathcal L_\psi=\frac{K_Q}{2}\mathcal Q^2-A\mathcal Y^{3/2}
 -\frac{\gamma}{2M_*^2}(\Delta_U\psi)^2.
 \label{eq:stage-a-force}
\end{equation}
The symbolic Stage-A calculation gives the necessary frozen-metric conditions
\begin{equation}
 c_{14}>0,\qquad c_1>0,\qquad c_{123}>0,
 \qquad K_Q>0,\quad A>0,\quad\gamma>0,
 \label{eq:stage-a-signs}
\end{equation}
and at zero background gradient the regulator produces
\begin{equation}
 \omega^2=\frac{\gamma}{K_QM_*^2}k^4.
 \label{eq:k4-dispersion}
\end{equation}
This is a pass only in the declared flat, constant-frame decoupling limit.

\subsection{Stage B frame-sector diagnostic}

A verified convention map to the Einstein--aether review
\cite{elingjacobsonmattingly2004} maps all four $c_i$ identically.  Substitution
of that source's coupled metric--aether result gives
\begin{align}
 s_2^2&=\frac{1}{1-c_{13}},\\
 s_1^2&=\frac{c_1-\tfrac12c_1^2+\tfrac12c_3^2}
 {c_{14}(1-c_{13})},\\
 s_0^2&=\frac{c_{123}(2-c_{14})}
 {c_{14}(1-c_{13})(2+c_{13}+3c_2)}.
 \label{eq:aether-speeds}
\end{align}
When all $c_i$ are scaled uniformly to zero, the spin-1 and spin-0 expressions
reduce to the Stage-A ratios $c_1/c_{14}$ and $c_{123}/c_{14}$.  This confirms
that those ratios are weak-metric-coupling limits, not the full mode speeds.
It is a checked literature substitution, not an independent ITSM-side
linearized Einstein-equation derivation, and it excludes the coupled
$\Phi$--$U$--$\psi$ system.

\subsection{Zero-gradient force-block factorization}

On the declared background $\bar\psi=\text{constant}$, every derivative of
the force field begins at first perturbative order.  Consequently,
\begin{equation}
 \mathcal Q^2=\dot\pi^2+\mathcal O(3),\qquad
 \mathcal Y^{3/2}=\mathcal O(3),\qquad
 (\Delta_U\psi)^2=(\nabla^2\pi)^2+\mathcal O(3),
\end{equation}
where metric and frame corrections enter only in the indicated higher-order
terms.  The force scalar therefore factorizes from the lapse, shift, frame and
condensate blocks at quadratic order for the declared truncation:
\begin{equation}
 S_\pi^{(2)}=\frac12\int\dd^4x\left[
 K_Q\dot\pi^2-\frac{\gamma}{M_*^2}(\nabla^2\pi)^2\right].
 \label{eq:zero-gradient-force-block}
\end{equation}
It carries one nonnegative $z=2$ scalar mode for $K_Q,\gamma>0$.  This bounded
result does not reduce the remaining metric--frame--condensate block.

A constant field redefinition $\psi_c=s\psi$ sends
\begin{equation}
 K_Q\to\frac{K_Q}{s^2},\quad
 A\to\frac{A}{s^3},\quad
 \gamma\to\frac{\gamma}{s^2},\quad
 C_m\to\frac{C_m}{s},\quad q\to sq.
\end{equation}
Thus $\gamma/K_Q$, $A/K_Q^{3/2}$, $C_m/\sqrt{K_Q}$ and $Aq/K_Q$ are invariant,
whereas $K_Q$ alone is not.  UVIR-003 can determine a structural stability
domain in invariant ratios, but a parent microscopic calculation or MAT-001
must fix the physical normalization and select a point in that domain.  The
earlier $K_Q\sim M_P^2$ estimate remains a speculative dimensional analogy.

\subsection{Causality and cutoff status}

Around a nonzero spatial background gradient of magnitude $q$, the force
fluctuation has
\begin{equation}
 \omega^2=\frac{3Aq(1+\cos^2\theta)}{K_Q}k^2
 +\frac{\gamma}{K_QM_*^2}k^4.
 \label{eq:force-dispersion}
\end{equation}
Thus the long-wavelength metric-light-cone threshold depends on the unmatched
ratio $Aq/K_Q$, while phase and group velocities grow without bound at large
$k$ inside the truncated dispersion.  Preferred-frame metric superluminality
is not by itself a proof of acausality, but the foliation-causality test and the
physical EFT domain have not been established.

Expansion of the cubic vertex supplies only the restricted, time-normalized
longitudinal IR NDA scale
\begin{equation}
 \Lambda_{\rm NDA,long}^{\rm (IR)}\sim\frac{K_Q^{3/4}}{\sqrt A}.
 \label{eq:long-nda}
\end{equation}
It is not the physical cutoff $\Lambda_{\rm EFT}$: transverse normalization,
$k^4$-dominated power counting and a loop or unitarity analysis remain open.
The comparison of the light-cone crossing scale with $\Lambda_{\rm EFT}$
therefore cannot yet be made.

A separate dimensional analogy, $K_Q=k_QM_P^2$, would place the
long-wavelength threshold at order $a_0$ for order-unity coefficients.  That
exercise is explicitly speculative: it is introduced only as a priority
signal for deriving $K_Q$, not as a property of ITSM.

\subsection{Conditional infrared operator}

Define the normalized spatial invariant
\begin{equation}
 Y=\frac{h^{\mu\nu}\nabla_\mu\psi\nabla_\nu\psi}{a_0^2}.
 \label{eq:Y}
\end{equation}
The target low-acceleration operator remains
\begin{equation}
 \mathcal L_{\psi,\rm IR}
 =-\frac{2C_{\rm IR}}{3}\mpl^2a_0^2Y^{3/2}
 +\mathcal O(Y^2,\nabla^2/\Lambda_{\rm EFT}^2).
 \label{eq:ir-action}
\end{equation}
Neither $C_{\rm IR}$, $K_Q$ nor the normalized matter vertex has been derived.
The preceding analytic Born--Infeld expression also remains rejected as the
source of this branch because its small-invariant equation is linear at
leading order.  UVIR-003 is therefore in progress and MAT-001 remains blocked.
